% ==============================================================
\section{Related Work}\label{sec:related}
% ==============================================================

This work addresses the problem of discovering
dense itemset regions from spatio-temporal transaction data.
This problem lies at the intersection of two research lineages:
(1)~the \emph{pattern mining lineage} aimed at enumerating itemsets, and
(2)~the \emph{dense block detection lineage} aimed at detecting
dense blocks in multi-dimensional spaces.
Below we organize both lineages and clarify
why this work is positioned at their intersection.

% ==============================================================
\subsection{Pattern Mining Lineage}
% ==============================================================

\paragraph{Frequent Itemset Mining (FIM).}
The Apriori algorithm by Agrawal et al.~\cite{agrawal1994apriori}
exploits the monotonicity of itemsets
(subsets always have no less support than their supersets)
to drastically prune the candidate space.
FP-Growth by Han et al.~\cite{han2000mining}
further accelerates mining using prefix trees.
These methods use global support over the entire
static transaction database as their criterion,
and thus suffer from the fundamental problem
of missing locally dense patterns~\cite{aprioriwindow}.

\paragraph{Local dense interval detection methods.}
Several methods based on \emph{intervals} between consecutive
occurrence timestamps have been proposed
for detecting locally dense intervals.
LPFIM~\cite{mahanta2005lpfim} splits intervals
when the gap between two consecutive occurrences exceeds
a threshold minthd1, and outputs intervals
whose local support rate exceeds a threshold
and whose length exceeds minthd2.
LPPM~\cite{fournierviger2021lppm} introduces
spillover accumulation for noise-robust interval detection.

These methods estimate density by evaluating
``the time difference between two consecutive occurrences
(inter-occurrence interval).''
However, as pointed out in~\cite{aprioriwindow},
since the gap parameter governs both pattern identification accuracy
and interval detection accuracy,
parameter settings that simultaneously keep both high
generally do not exist.
Moreover, these methods are \emph{inherently non-extensible to
multiple dimensions}.
Occurrence points on a $d$-dimensional grid
have no linear ordering,
making ``the interval between two consecutive occurrences''
ill-defined.

Apriori-Window~\cite{aprioriwindow} is an exact method
that resolves this structural limitation.
By \emph{directly} computing the number of occurrences
within a sliding window $[l, l+W]$
rather than via inter-occurrence intervals,
it evaluates local frequency without requiring a gap parameter.
Combined with anti-monotonicity-based Apriori pruning
and $t$-axis stride, it exhaustively and exactly
enumerates dense intervals.
Comparative experiments demonstrated that
Apriori-Window significantly outperforms
LPFIM and LPPM~\cite{aprioriwindow}.

\emph{Why multi-dimensionalization is possible}:
The definition ``number of occurrences within a rectangular window''
naturally extends to $d$ dimensions
(simply replacing the window with a $d$-dimensional hyperrectangle
$\prod_k [s_k, s_k+W_k]$).
In contrast to interval-based methods
that depend on linear ordering,
the design of Apriori-Window does not require linear ordering.
This work builds on this observation to extend
Apriori-Window from one dimension to arbitrary $d$ dimensions.

\paragraph{Colocation mining.}
Colocation mining~\cite{shekhar2001colocation}
evaluates patterns where multiple feature types
co-exist in geographic proximity
using the Participation Index (PI).
Qian et al.~\cite{qian2009stcolocation} extended this
to weighted sliding windows (STCP-Miner).
These methods deal with spatial proximity-based co-occurrence
rather than transaction-based co-occurrence
and do not extract dense regions.
However, they can be compared in terms of
pattern identification accuracy (F1) (\S\ref{sec:ex1_1}).

% ==============================================================
\subsection{Region Extraction Lineage}
% ==============================================================

% --- Dense block detection ---
\paragraph{Dense block detection methods.}
Dense block detection methods
including M-Zoom~\cite{shin2016mzoom}
and others (tensor:~\cite{jiang2015crossspot,shin2018mbiz,shin2017dcube,shin2017densealert,feng2020catchcore};
graph: SpecGreedy~\cite{feng2020specgreedy})
detect dense blocks in multi-dimensional data.
They select value sets $B_k \subseteq D_k$ along each dimension
and perform greedy search to maximize
the mass-to-volume ratio of the Cartesian-product block
$\prod_k B_k$.

These methods have two structural limitations.

\emph{First, output is restricted to Cartesian-product blocks.}
For irregularly-shaped dense regions,
they inevitably include excessive areas
or miss dense cells,
resulting in low Jaccard even when GT patterns are provided
(demonstrated in \S\ref{sec:ex1_1}).

\emph{Second, pattern enumeration is inaccurate.}
While adding items as an additional dimension
allows enumeration to be attempted,
the Cartesian-product constraint causes all subsets
of the itemset $C_B$ within a block
to be assumed dense in the same region,
generating false positives (demonstrated in \S\ref{sec:ex1_1}).

% --- Spatial statistics ---
\paragraph{Spatial statistics.}
LISA~\cite{anselin1995lisa} identifies spatial hotspots
using local Moran's $I$, and
Getis-Ord $G_i^*$~\cite{ord1995getis} evaluates
local concentration via $z$-scores.
Both target \emph{spatial dimensions only};
running them per time slice
loses temporal continuity.

% --- Spatial/spatio-temporal scan statistics ---
\paragraph{Spatial and spatio-temporal scan statistics.}
Kulldorff's~\cite{kulldorff1997spatial,kulldorff2001prospective}
spatial/spatio-temporal scan statistics detect
significant clusters using circular/cylindrical scanning windows.
FlexScan~\cite{tango2005flexscan} extends this
to irregularly-shaped windows.

All above statistical methods are designed
to detect per-variable anomalous concentrations
and cannot enumerate co-occurrence patterns (itemsets)
within transactions.
Enumerating all patterns would require
$O(2^{|\mathcal{I}|})$ independent runs,
and integration with Apriori pruning is not part of their design.

% --- Spatio-temporal clustering ---
\paragraph{Spatio-temporal density-based clustering (ST-DBSCAN, etc.).}
ST-DBSCAN~\cite{birant2007stdbscan} extends
DBSCAN~\cite{ester1996dbscan} to spatio-temporal data
with independent thresholds $\varepsilon_1, \varepsilon_2$
for spatial and temporal distances.
While it outputs arbitrarily-shaped clusters,
first, it defines density by point-neighborhood counts;
when occurrences are spatially uniform,
the entire dataset becomes a single cluster,
failing to identify localized concentrations.
Second, point-neighborhood density lacks anti-monotonicity,
making independent processing of all patterns
$O(2^{|\mathcal{I}|})$ unavoidable.

All above region extraction methods lack the concept of itemsets,
but they can be evaluated as region extractors
when GT patterns and their occurrence densities are provided
(\S\ref{sec:ex1_1}).

% ==============================================================
\subsection{Positioning of This Work}
% ==============================================================

The above analysis shows that no existing method
simultaneously enumerates patterns
and extracts multi-dimensional dense regions.
This work achieves both within a single framework
by exploiting the anti-monotonicity
of windowed-occurrence-based density
to naturally integrate Apriori pruning.
Table~\ref{tab:related} illustrates this positioning.

\paragraph{Why no stronger baseline exists.}
The gap in Table~\ref{tab:related} reflects a \emph{structural} incompatibility:
pattern mining methods cannot output spatio-temporal windows,
and region extraction methods lose Apriori pruning when applied independently per itemset.
Faster variants within the same categories share these limitations;
the chosen baselines therefore represent each distinct methodological category.

\begin{table}[t]
\centering
\caption{Comparison with related methods.
``Itemset'' refers to the ability to enumerate patterns;
``Dense region'' refers to the ability to output
spatio-temporally localized dense regions.}
\label{tab:related}
\small
\begin{tabular}{lccc}
\toprule
Method & \shortstack{Item-\\set} & \shortstack{ST dense\\region} & \shortstack{Region\\shape} \\
\midrule
Apriori~\cite{agrawal1994apriori}
                                          & $\checkmark$ & $\times$     & ---          \\
LPFIM~\cite{mahanta2005lpfim} et al.
                                          & $\checkmark$ & 1D only      & Interval     \\
Apriori-Window~\cite{aprioriwindow}       & $\checkmark$ & 1D only      & Interval     \\
Colocation~\cite{shekhar2001colocation}
                                          & $\checkmark$ & $\times$     & ---          \\
STCP-Miner~\cite{qian2009stcolocation}   & $\checkmark$ & $\times$     & ---          \\
\midrule
M-Zoom~\cite{shin2016mzoom} et al.       & $\triangle$$^*$     & $\checkmark$ & Cartesian only \\
SaTScan~\cite{kulldorff1997spatial}       & $\times$     & $\checkmark$ & Circle/cylinder \\
FlexScan~\cite{tango2005flexscan}        & $\times$     & $\checkmark$ & Irregular    \\
LISA~\cite{anselin1995lisa}
  / $G_i^*$~\cite{ord1995getis}          & $\times$     & Spatial only & Arbitrary    \\
ST-DBSCAN~\cite{birant2007stdbscan} et al. & $\times$   & $\checkmark$ & Arbitrary    \\
\midrule
\textbf{Ours}                              & $\checkmark$ & $\checkmark$ & Arbitrary    \\
\bottomrule
\multicolumn{4}{l}{\scriptsize $^*$Patterns can be output by including items as a dimension,} \\
\multicolumn{4}{l}{\scriptsize \hspace{1em}but Cartesian constraints cause false positives and preclude region separation.}
\end{tabular}
\end{table}
